\section{Sensing: The Accelerometer \& Data Acquisition}
\label{sec:sensing}

Everything Cadence knows about a wearer's gait arrives through a single channel: the
motion of one ankle, measured many times a second. Before any spectral analysis,
freeze index, or neural network can run, that motion must be turned into clean, evenly
spaced numbers. This section explains the sensor that does it --- the LIS3DH MEMS
accelerometer on the Circuit Playground Express (CPX) --- and the disciplined,
drift-free loop on the SAMD21 microcontroller that samples it. The choices here
(which axes, what range, how fast, what units) set the ceiling on everything
downstream, so they deserve to be understood from first principles.

\subsection{What a MEMS accelerometer actually measures}

A MEMS (Micro-Electro-Mechanical System) accelerometer is, quite literally, a
microscopic spring--mass system etched into silicon. A tiny proof mass is suspended
on flexible silicon beams that behave as springs. When the chip accelerates, the mass
lags behind by inertia and the beams flex; that deflection is read out
\emph{capacitively}. The moving mass carries interleaved fingers that sit between
fixed fingers, forming a pair of capacitors. As the mass shifts by a few nanometres,
one gap narrows and the other widens, the capacitances change in opposite directions,
and the chip's analogue front end converts that differential capacitance into a
digital number proportional to acceleration along that axis.

\begin{intuition}
Picture a ball on a spring inside a sealed box. Shake the box and the ball swings;
how far it swings tells you how hard you shook. A MEMS accelerometer is that toy
shrunk onto a grain of silicon, with the ``how far'' read not by eye but by measuring
a capacitance. Crucially, it cannot tell ``being pushed'' from ``resting in
gravity'': sitting still on a table, the proof mass still sags under its own weight,
so the sensor always reports a steady $\approx\SI{1}{g}$ pull. That ever-present
offset is a feature we exploit later, not a fault.
\end{intuition}

The LIS3DH integrates three such structures at right angles, giving acceleration
along three orthogonal axes $a_x, a_y, a_z$ simultaneously. We mount it in an ankle
sleeve, so the axes are fixed to the leg rather than to the world. Figure~\ref{fig:axes}
shows the convention: one axis runs roughly vertically (and so carries most of
gravity when standing), one along the direction of travel, and one across the body.
Because a sleeve can be donned at slightly different angles each day, we never trust
any single axis on its own --- the orientation is not guaranteed. The remedy is the
magnitude $a=\sqrt{a_x^2+a_y^2+a_z^2}$, introduced in the next section, which is
identical no matter how the device is rotated.

% ---------------------------------------------------------------------------
\begin{figure}[H]
\centering
\resizebox{0.95\textwidth}{!}{%
\begin{tikzpicture}[font=\small,>=Latex]

  % ---- LEFT: leg + ankle axes ----
  \begin{scope}[shift={(0,0)}]
    % stylised lower leg (kept narrow + to the left so the axis triad has clear space)
    \fill[mist] (-0.45,0) -- (0.45,0) -- (0.32,4.2) -- (-0.32,4.2) -- cycle;
    \draw[rule,thick] (-0.45,0) -- (0.45,0) -- (0.32,4.2) -- (-0.32,4.2) -- cycle;
    % foot
    \fill[mist] (-0.45,0) -- (1.55,0) -- (1.55,-0.45) -- (-0.45,-0.45) -- cycle;
    \draw[rule,thick] (-0.45,0) -- (1.55,0) -- (1.55,-0.45) -- (-0.45,-0.45) -- cycle;
    \node[mute,font=\footnotesize] at (0.9,-0.74) {foot};
    \node[mute,font=\footnotesize,rotate=90] at (-0.82,2.35) {lower leg};

    % sleeve band
    \fill[paleteal,opacity=0.9] (-0.42,0.7) rectangle (0.40,1.65);
    \draw[teal,thick] (-0.42,0.7) rectangle (0.40,1.65);
    % sensor board
    \fill[teal] (-0.18,0.99) rectangle (0.18,1.35);
    \node[white,font=\scriptsize\bfseries] at (0,1.17) {CPX};
    % callout label to the lower-left with a thin leader (clear of the triad)
    \node[teal,font=\footnotesize,anchor=east] at (-1.05,0.55) {ankle sleeve};
    \draw[teal,thin] (-0.95,0.6) -- (-0.42,0.95);

    % axis triad anchored at the board, opening to the RIGHT into empty space
    \coordinate (O) at (0,1.17);
    \draw[->,teal,very thick]  (O) -- ++(0,1.75)    node[above,teal,font=\footnotesize]{$z$ (vertical, gravity)};
    \draw[->,steel,very thick] (O) -- ++(2.0,0)     node[right,steel,font=\footnotesize]{$x$ (forward)};
    \draw[->,navy,very thick]  (O) -- ++(1.0,-0.78) node[below right,navy,font=\footnotesize]{$y$ (lateral)};
  \end{scope}

  % ---- RIGHT: capacitive proof-mass schematic ----
  \begin{scope}[shift={(6.6,2.0)}]
    \node[navy,font=\bfseries] at (1.6,2.35) {Inside one axis: a capacitive spring--mass};
    % anchors
    \fill[mute] (-0.2,1.6) rectangle (0.2,1.85);
    \fill[mute] (3.0,1.6) rectangle (3.4,1.85);
    \fill[mute] (-0.2,-1.85) rectangle (0.2,-1.6);
    \fill[mute] (3.0,-1.85) rectangle (3.4,-1.6);
    % springs (zig-zag)
    \draw[ink,thick] (0,1.6) -- (0.25,1.3) -- (-0.05,1.0) -- (0.25,0.7) -- (0.1,0.45);
    \draw[ink,thick] (3.2,1.6) -- (2.95,1.3) -- (3.25,1.0) -- (2.95,0.7) -- (3.1,0.45);
    \draw[ink,thick] (0,-1.6) -- (0.25,-1.3) -- (-0.05,-1.0) -- (0.25,-0.7) -- (0.1,-0.45);
    \draw[ink,thick] (3.2,-1.6) -- (2.95,-1.3) -- (3.25,-1.0) -- (2.95,-0.7) -- (3.1,-0.45);
    % proof mass (movable)
    \fill[palenavy] (0.1,-0.45) rectangle (3.1,0.45);
    \draw[navy,very thick] (0.1,-0.45) rectangle (3.1,0.45);
    \node[navy,font=\scriptsize\bfseries] at (1.6,0) {proof mass};
    % moving fingers
    \foreach \x in {0.5,1.1,1.7,2.3,2.9}{
      \fill[navy] (\x,0.45) rectangle (\x+0.10,0.95);
      \fill[navy] (\x,-0.95) rectangle (\x+0.10,-0.45);
    }
    % fixed fingers (offset) -- top = C1, bottom = C2
    \foreach \x in {0.8,1.4,2.0,2.6}{
      \fill[teal] (\x,0.55) rectangle (\x+0.10,1.05);
      \fill[freeze] (\x,-1.05) rectangle (\x+0.10,-0.55);
    }
    \node[teal,right] at (3.25,0.8) {\footnotesize $C_1$};
    \node[freeze,right] at (3.25,-0.8) {\footnotesize $C_2$};
    % acceleration arrow
    \draw[->,amber,very thick] (1.6,1.45) -- (1.6,1.15);
    \node[amber,right] at (1.75,1.3) {\footnotesize mass lags $\Rightarrow$ $C_1\!\neq\!C_2$};
    \node[mute,align=center] at (1.6,-2.25)
      {\footnotesize deflection of a few nm $\to$ differential capacitance $\to$ digital count};
  \end{scope}
\end{tikzpicture}}
\caption{Left: the LIS3DH is worn in an ankle sleeve, fixing its three orthogonal
axes to the leg --- $z$ vertical (and so loaded by gravity at rest), $x$ forward,
$y$ lateral. Right: along each axis a silicon proof mass on spring beams deflects
under acceleration, unbalancing a pair of capacitors $C_1, C_2$; that differential
capacitance is digitised. Because the sleeve's orientation is not guaranteed day to
day, Cadence never trusts a single axis and instead works from the rotation-invariant
magnitude.}
\label{fig:axes}
\end{figure}

\subsection{Range: why \texorpdfstring{$\pm\SI{8}{g}$}{+/-8 g} and not \texorpdfstring{$\pm\SI{2}{g}$}{+/-2 g}}

The LIS3DH can be configured for several full-scale ranges. A smaller range gives
finer resolution per count, which is tempting --- but a range that is too small
\emph{clips}: any acceleration beyond the limit is recorded as a flat-topped maximum,
destroying exactly the peaks we care about. Walking is not gentle. Each
\emph{heel strike} sends a sharp impulse up the leg that comfortably exceeds
$\pm\SI{2}{g}$. If we ran at $\pm\SI{2}{g}$, those impacts would saturate the sensor,
flattening the very spikes that distinguish vigorous gait from a freeze. We therefore
run the LIS3DH at $\pm\SI{8}{g}$, which leaves clean headroom above every heel strike.

\begin{takeaway}
A heel strike exceeds $\pm\SI{2}{g}$, so the small range would clip. Cadence runs at
$\pm\SI{8}{g}$ --- never trading dynamic range for resolution when the peaks are the
signal.
\end{takeaway}

Figure~\ref{fig:capture} is real data from the worksheet capture (\code{capture1.csv},
\SI{64}{Hz}, \SI{41.4}{s}). The top trace is the acceleration magnitude in milli-$g$.
Two facts jump out and tie directly to the choices in this section. First, the
baseline sits near $\SI{1000}{mg}$ even when the wearer is still --- that is the
$\approx\SI{1}{g}$ of gravity the proof mass always feels, the DC offset we remove
before spectral analysis. Second, during the freeze phase the magnitude swings to
roughly $\SI{3500}{mg}$, far above $\SI{2000}{mg}$: had the range been $\pm\SI{2}{g}$
those excursions would have been sheared flat at the rail.

\begin{figure}[H]
\centering
\includegraphics[width=0.92\linewidth]{capture1_trace}
\caption{Real \SI{64}{Hz} capture (\code{capture1.csv}, \SI{41.4}{s}). Top: the
acceleration magnitude $\norm{a}$ in milli-$g$, sitting on a $\approx\SI{1000}{mg}$
gravity baseline; the freeze phase drives it well past $\SI{2000}{mg}$, which a
$\pm\SI{2}{g}$ range would clip. Bottom: the corresponding Freeze Index, computed
downstream from this same magnitude (Section~\ref{sec:sensing} feeds it). The sensor's
job is to deliver these excursions \emph{without saturating}; the analysis that turns
them into a verdict comes later.}
\label{fig:capture}
\end{figure}

\subsection{Sampling rate, Nyquist, and why \texorpdfstring{\SI{64}{Hz}}{64 Hz} is plenty}

Cadence samples all three axes at $f_s=\SI{64}{Hz}$ --- 64 fresh readings per second.
The right way to judge whether that is fast enough is the Nyquist--Shannon sampling
theorem: a sampled signal can faithfully represent frequencies only up to half the
sampling rate. Here the Nyquist frequency is $f_s/2 = \SI{32}{Hz}$. Every frequency
band Cadence cares about sits far below that ceiling: the locomotor (walking-rhythm)
band at \SIrange{0.5}{3}{Hz}, the freeze band at \SIrange{3}{8}{Hz}, and the
rest-tremor band at \SIrange{4}{6}{Hz}. The highest of these tops out at \SI{8}{Hz},
four times below Nyquist (Figure~\ref{fig:nyquist}). There is comfortable margin, and
no aliasing of the freeze signature into a misleading frequency.

\begin{mathsbox}
\textbf{Nyquist, in one line.} To capture a frequency $f$ without aliasing you must
sample faster than twice it:
\begin{equation}
f_s > 2 f_{\max}
\qquad\Longleftrightarrow\qquad
f_{\max} < \frac{f_s}{2} \;=\; \frac{\SI{64}{Hz}}{2} \;=\; \SI{32}{Hz}.
\label{eq:nyquist}
\end{equation}
Our highest band of interest is \SI{8}{Hz}. Since $8 < 32$ with a factor of four to
spare, \SI{64}{Hz} represents the \SIrange{3}{8}{Hz} freeze trembling exactly.
Sample \emph{too slowly} --- say below \SI{16}{Hz} --- and an \SI{8}{Hz} tremor would
fold down (\emph{alias}) to a lower, fictitious frequency and corrupt the Freeze
Index. \SI{64}{Hz} avoids that entirely while keeping data rates and on-board compute
modest.
\end{mathsbox}

\begin{figure}[H]
\centering
\resizebox{0.95\textwidth}{!}{%
\begin{tikzpicture}[font=\small]
\begin{axis}[
    width=14cm, height=5.2cm,
    xmin=0, xmax=34, ymin=0, ymax=1.25,
    xlabel={frequency (Hz)}, ylabel={},
    ytick=\empty,
    axis y line=none,
    axis x line=bottom,
    xtick={0,3,6,8,16,32},
    xticklabel style={font=\footnotesize},
    every axis x label/.style={at={(axis description cs:0.5,-0.16)},anchor=north},
    clip=false,
]
  % locomotor band 0.5-3
  \addplot[draw=none,fill=steel,fill opacity=0.30] coordinates {(0.5,0)(3,0)(3,1)(0.5,1)} \closedcycle;
  % freeze band 3-8
  \addplot[draw=none,fill=freeze,fill opacity=0.28] coordinates {(3,0)(8,0)(8,1)(3,1)} \closedcycle;
  % rest-tremor band 4-6 (overlaps freeze) -- hatch
  \addplot[draw=none,fill=purple,fill opacity=0.22] coordinates {(4,0)(6,0)(6,0.62)(4,0.62)} \closedcycle;

  % band labels -- each centred over its OWN band on a separate baseline so they never collide
  \node[steel,font=\footnotesize\bfseries,align=center,anchor=south] at (axis cs:1.75,1.30) {locomotor\\\SIrange{0.5}{3}{Hz}};
  \node[freeze,font=\footnotesize\bfseries,align=center,anchor=south] at (axis cs:5.5,1.00) {freeze\\\SIrange{3}{8}{Hz}};
  \node[purple,font=\footnotesize,align=center] at (axis cs:5.0,0.42) {rest\\tremor};

  % Nyquist line
  \draw[navy,very thick,dashed] (axis cs:32,0) -- (axis cs:32,1.15);
  \node[navy,font=\footnotesize\bfseries,align=center,anchor=south] at (axis cs:32,1.16)
     {Nyquist\\$f_s/2=\SI{32}{Hz}$};

  % headroom brace from 8 to 32
  \draw[<->,amber,thick] (axis cs:8,0.30) -- (axis cs:32,0.30);
  \node[amber,font=\footnotesize] at (axis cs:20,0.40) {comfortable margin};
\end{axis}
\end{tikzpicture}}
\caption{Where the action is, against the sampling ceiling. All three bands Cadence
analyses lie below \SI{8}{Hz}; the Nyquist frequency for \SI{64}{Hz} sampling is
\SI{32}{Hz}, leaving a four-fold margin so nothing aliases. (The rest-tremor band
\SIrange{4}{6}{Hz} sits inside the freeze band.)}
\label{fig:nyquist}
\end{figure}

\subsection{Units and the gravity offset}

On the device the SAMD21 works in SI units, $\si{m/s^2}$, because that is what the
LIS3DH driver returns. On the host side we convert to milli-$g$ ($\SI{1}{g}\approx
\SI{9.81}{m/s^2}$, so $\SI{1}{mg}\approx\SI{9.81e-3}{m/s^2}$) because milli-$g$ is the
natural language of human movement: a still limb reads $\approx\SI{1000}{mg}$, a brisk
heel strike a few thousand. Whichever unit is used, one preprocessing step is
non-negotiable. The proof mass always feels gravity, so every axis carries a constant
$\approx\SI{1}{g}$ bias --- a strong DC component. Before any spectral analysis we
subtract the mean of the magnitude window, removing that DC term so the power spectrum
reflects \emph{movement}, not the constant pull of the Earth.

\begin{keyidea}
Gravity is a DC offset, not motion. Subtracting the window mean from the magnitude
zeroes the $\approx\SI{1000}{mg}$ baseline, so the spectrum that feeds the Freeze
Index measures how the leg is \emph{shaking}, not the fact that it is on a planet.
Because the Freeze Index is later formed as a ratio of band powers, the choice of
unit (mg versus g) cancels --- but the DC removal does not, and must happen first.
\end{keyidea}

\subsection{Drift-free acquisition on the SAMD21}

A spectral analysis is only as honest as its time axis. If samples arrive at uneven
intervals --- because the firmware did some I2C work, drew the OLED, or updated the
NeoPixels in between --- the assumed \SI{64}{Hz} grid is a lie, and a \SI{6}{Hz}
tremor can smear into the wrong bin. The classic beginner's mistake is to sample,
then \code{delay(15)} (roughly $1/64\,\si{s}$): the delay does not account for the
time the sampling and processing themselves took, so the true period creeps and
\emph{drifts}.

Cadence instead uses fixed-rate, non-blocking scheduling. It records the timestamp of
the next sample that is \emph{due} and advances that target by a fixed period each
time, rather than sleeping for a fixed period. The loop never blocks: between samples
it is free to run the on-board detector, refresh the OLED, and drive the LED ring,
yet a sample is taken the instant its scheduled deadline passes. Errors do not
accumulate, because each deadline is computed from the previous \emph{target}, not
from when the previous sample actually fired.

\begin{lstlisting}[style=cpp]
// Fixed-rate, non-blocking 64 Hz sampler on the SAMD21 (sketch)
const uint32_t PERIOD_US = 1000000UL / 64;   // 15625 us between samples
uint32_t next_due = micros();                // first deadline = now

void loop() {
  if ((int32_t)(micros() - next_due) >= 0) { // deadline passed?
    next_due += PERIOD_US;                    // advance by a FIXED step
    float ax, ay, az;
    lis3dh_read(&ax, &ay, &az);               // m/s^2, +/-8 g range
    push_sample(ax, ay, az);                  // into the ring buffer
  }
  service_detector();   // FI / gate / debounce  -- never blocks
  service_display();    // OLED + NeoPixel ring  -- never blocks
}
\end{lstlisting}

\begin{intuition}
Think of it as catching a train by the clock, not by a stopwatch. ``Sleep 15 ms''
restarts the stopwatch \emph{after} you have already spent time reading the sensor and
drawing the screen, so each cycle runs a little long and the schedule slips. ``Take a
sample whenever the wall clock crosses the next 15.625 ms tick'' keeps every sample
pinned to an absolute deadline --- the timing never drifts, no matter how much work
happens in between.
\end{intuition}

Because the advance is by a fixed \code{PERIOD\_US} computed from the target rather
than the actual fire time, the long-run average rate stays locked to \SI{64}{Hz} and
the time base under the Welch analysis is trustworthy. The cast to \code{int32\_t} in
the comparison makes the deadline test robust to the microsecond counter wrapping
around. Two firmware roles share this same acquisition core: the
\emph{recorder} (\code{cpx\_fog\_logger.ino}) streams samples to the CPX's
\SI{2}{MB} flash for later analysis, and the \emph{monitor}
(\code{cpx\_fog\_standalone.ino}) feeds them straight into the live detector and
display. The sensing front end is identical in both; only what happens to the samples
afterwards differs.

\begin{takeaway}
One ankle accelerometer, sampled at a rock-steady \SI{64}{Hz} over $\pm\SI{8}{g}$,
delivers gravity-baselined, alias-free motion. Get the range, rate, and timing right
here and every downstream stage --- magnitude, spectrum, Freeze Index, dashboard ---
inherits clean data; get them wrong and no clever analysis can recover the lost peaks.
\end{takeaway}
